% answers.tex —— 2014 年·北京卷·文科 · 参考答案与详细解析片段
% 由 tools/generate_projects.py 依来源总账自 2014/2014_数学真题_北京文_试题.md + 2014/2014_数学真题_北京文_解析正文.md + 2014/2014_数学真题_北京文_答案要点.md 生成，勿手改（改源 MD 后重跑）


\begin{answer}
\gksolhead{1}{}

\textbf{题目}：

若集合 $A=\{0,1,2,4\}$，$B=\{1,2,3\}$，则 $A\cap B=$\mbox{（　　）}

（A）$\{0,1,2,3,4\}$；（B）$\{0,4\}$；（C）$\{1,2\}$；（D）$\{3\}$。

\Needspace{4\baselineskip}
\textbf{解析}

\Needspace{5\baselineskip}
$A\cap B$ 由同时属于 $A$ 与 $B$ 的元素组成。逐个检验：$0\in A$ 但 $0\notin B$；$1$，$2$ 同时属于 $A$，$B$；$4\in A$ 但 $4\notin B$；$3\in B$ 但 $3\notin A$。故
\nopagebreak[4]
\[A\cap B=\{1,2\},\]
选 C。

\gkanswer{答案：C。$A$，$B$ 的公共元素为 $1$，$2$，故 $A\cap B=\{1,2\}$。}

\end{answer}
\begin{answer}
\gksolhead{2}{}

\textbf{题目}：

下列函数中，定义域是 $\mathbf R$ 且为增函数的是\mbox{（　　）}

（A）$y=\mathrm e^{-x}$；（B）$y=x^3$；（C）$y=\ln x$；（D）$y=|x|$。

\Needspace{4\baselineskip}
\textbf{解析}

逐项检查定义域与单调性：

\begin{itemize}
\item $y=\mathrm e^{-x}$ 的定义域为 $\mathbf R$，但它是减函数；
\item $y=x^3$ 的定义域为 $\mathbf R$，且对任意 $x_1<x_2$ 有 $x_1^3<x_2^3$，是增函数；
\item $y=\ln x$ 的定义域为 $(0,+\infty)$，不是 $\mathbf R$；
\item $y=|x|$ 的定义域为 $\mathbf R$，但在 $(-\infty,0)$ 上递减，不是 $\mathbf R$ 上的增函数。
\end{itemize}

故选 B。

\gkanswer{答案：B。$y=x^3$ 定义域为 $\mathbf R$ 且严格递增；其余三项或为减函数，或定义域不是 $\mathbf R$，或不单调。}

\end{answer}
\begin{answer}
\gksolhead{3}{}

\textbf{题目}：

已知向量 $\boldsymbol a=(2,4)$，$\boldsymbol b=(-1,1)$，则 $2\boldsymbol a-\boldsymbol b=$\mbox{（　　）}

（A）$(5,7)$；（B）$(5,9)$；（C）$(3,7)$；（D）$(3,9)$。

\Needspace{4\baselineskip}
\textbf{解析}

\Needspace{5\baselineskip}
由数乘与减法的坐标运算，
\nopagebreak[4]
\[2\boldsymbol a-\boldsymbol b=2(2,4)-(-1,1)=(4,8)-(-1,1)=(5,7),\]
选 A。

\gkanswer{答案：A。$2\boldsymbol a-\boldsymbol b=(4,8)-(-1,1)=(5,7)$。}

\end{answer}
\begin{answer}
\gksolhead{4}{}

\textbf{题目}：

执行如图所示的程序框图，则输出的 $S$ 值为\mbox{（　　）}
\nopagebreak[4]
\begin{center}% ============================================================
% f2014_beijing_arts_q04_flowchart.tex
% 2014 北京卷（文科）第 4 题：程序框图
% 来源 PDF：普通高考/2014/2014北京文.pdf
% 原页：第 1 页（A3），印面 863，左栏第 4 题题干下方
% 结构：开始 → k=0, S=0 → 判断 k<3；是则 S=S+2^k、k=k+1 并返回判断；
%       否则向右输出 S，再向下结束。
% ============================================================
\begin{tikzpicture}[>=Stealth, line width=0.7pt, font=\small]
  % 开始
  \draw[rounded corners=6pt] (-0.95,0.45) rectangle (0.95,-0.45);
  \node at (0,0) {开始};
  \draw[->] (0,-0.45) -- (0,-1.05);

  % 初始化
  \draw (-1.75,-1.05) rectangle (1.75,-2.05);
  \node at (0,-1.55) {$k=0$，$S=0$};
  \draw (0,-2.05) -- (0,-2.35);
  \draw[->] (0,-2.35) -- (0,-2.65);

  % 判断菱形
  \draw (0,-2.65) -- (1.55,-3.45) -- (0,-4.25) -- (-1.55,-3.45) -- cycle;
  \node at (0,-3.45) {$k<3$};

  % 是：循环体
  \draw[->] (0,-4.25) -- (0,-4.95);
  \node at (-0.22,-4.62) {是};
  \draw (-1.75,-4.95) rectangle (1.75,-5.95);
  \node at (0,-5.45) {$S=S+2^k$};
  \draw[->] (0,-5.95) -- (0,-6.35);
  \draw (-1.75,-6.35) rectangle (1.75,-7.35);
  \node at (0,-6.85) {$k=k+1$};

  % 回边：由 k=k+1 左侧绕回判断框之前的竖线
  \draw (-1.75,-6.85) -- (-2.85,-6.85) -- (-2.85,-2.35);
  \draw[->] (-2.85,-2.35) -- (0,-2.35);

  % 否：输出
  \draw[->] (1.55,-3.45) -- (3.35,-3.45);
  \node at (2.35,-3.20) {否};
  \draw[->] (3.35,-3.45) -- (3.35,-3.98);
  \draw (1.95,-4.68) -- (4.45,-4.68) -- (4.15,-3.98) -- (1.65,-3.98) -- cycle;
  \node at (3.05,-4.33) {输出 $S$};
  \draw[->] (3.05,-4.68) -- (3.05,-5.12);

  % 结束
  \draw[rounded corners=6pt] (2.15,-5.12) rectangle (3.95,-6.02);
  \node at (3.05,-5.57) {结束};
\end{tikzpicture}
\end{center}
\nopagebreak[4]
（A）$1$；（B）$3$；（C）$7$；（D）$15$。

\Needspace{4\baselineskip}
\textbf{解析}

初始 $k=0$，$S=0$。循环条件是 $k<3$，每轮先累加 $2^k$ 再让 $k$ 增加 $1$：

\begin{itemize}
\item $k=0<3$：$S=0+2^0=1$，$k=1$；
\item $k=1<3$：$S=1+2^1=3$，$k=2$；
\item $k=2<3$：$S=3+2^2=7$，$k=3$；
\item $k=3<3$ 不成立，输出 $S=7$。
\end{itemize}

即 $S=2^0+2^1+2^2=7$，选 C。

\gkanswer{答案：C。$S=2^0+2^1+2^2=7$。}

\end{answer}
\begin{answer}
\gksolhead{5}{}

\textbf{题目}：

设 $a$，$b$ 是实数，则“$a>b$”是“$a^2>b^2$”的\mbox{（　　）}

（A）充分且不必要条件；（B）必要且不充分条件；（C）充分必要条件；（D）既不充分也不必要条件。

\Needspace{4\baselineskip}
\textbf{解析}

取 $a=1$，$b=-2$，则 $a>b$ 成立，但 $a^2=1<4=b^2$，故“$a>b$”推不出“$a^2>b^2$”，不是充分条件。

取 $a=-2$，$b=1$，则 $a<b$，但 $a^2=4>1=b^2$，故“$a^2>b^2$”也推不出“$a>b$”，不是必要条件。

因此是既不充分也不必要条件，选 D。

\gkanswer{答案：D。$a=1$，$b=-2$ 反证不充分；$a=-2$，$b=1$ 反证不必要。}

\end{answer}
\begin{answer}
\gksolhead{6}{}

\textbf{题目}：

已知函数 $f(x)=\dfrac6x-\log_2x$，在下列区间中，包含 $f(x)$ 零点的区间是\mbox{（　　）}

（A）$(0,1)$；（B）$(1,2)$；（C）$(2,4)$；（D）$(4,+\infty)$。

\Needspace{4\baselineskip}
\textbf{解析}

\Needspace{5\baselineskip}
$f(x)$ 的定义域为 $(0,+\infty)$。其中 $\dfrac6x$ 与 $-\log_2x$ 都在 $(0,+\infty)$ 上单调递减，故 $f(x)$ 在 $(0,+\infty)$ 上单调递减，零点至多一个。计算端点值：
\nopagebreak[4]
\[f(2)=\frac62-\log_22=3-1=2>0,\qquad
f(4)=\frac64-\log_24=\frac32-2=-\frac12<0 .\]
由零点存在定理，零点在区间 $(2,4)$ 内，选 C。（又 $f(1)=6>0$，故 $(0,1)$，$(1,2)$ 内均无零点。）

\gkanswer{答案：C。$f(2)=2>0$，$f(4)=-\dfrac12<0$，且 $f$ 在 $(0,+\infty)$ 上递减。}

\end{answer}
\begin{answer}
\gksolhead{7}{}

\textbf{题目}：

已知圆 $C:(x-3)^2+(y-4)^2=1$ 和两点 $A(-m,0)$，$B(m,0)$（$m>0$），若圆 $C$ 上存在点 $P$，使得 $\angle APB=90^\circ$，则 $m$ 的最大值为\mbox{（　　）}

（A）$7$；（B）$6$；（C）$5$；（D）$4$。

\Needspace{4\baselineskip}
\textbf{解析}

\Needspace{5\baselineskip}
$\angle APB=90^\circ$ 等价于 $P$ 在以 $AB$ 为直径的圆上，即
\nopagebreak[4]
\[x^2+y^2=m^2 .\]
\Needspace{5\baselineskip}
于是问题化为该圆与圆 $C:(x-3)^2+(y-4)^2=1$ 有公共点。两圆圆心分别为 $(0,0)$，$(3,4)$，圆心距
\nopagebreak[4]
\[d=\sqrt{3^2+4^2}=5,\]
\Needspace{5\baselineskip}
半径分别为 $m$ 与 $1$，有公共点的充要条件是
\nopagebreak[4]
\[|m-1|\leqslant 5\leqslant m+1,\]
即 $m\leqslant 6$ 且 $m\geqslant 4$。所以 $4\leqslant m\leqslant 6$，$m$ 的最大值为 $6$，选 B。

\gkanswer{答案：B。$P$ 在以 $AB$ 为直径的圆 $x^2+y^2=m^2$ 上，两圆有公共点给出 $4\leqslant m\leqslant 6$。}

\end{answer}
\begin{answer}
\gksolhead{8}{}

\textbf{题目}：

加工爆米花时，爆开且不糊的粒数的百分比称为“可食用率”。在特定条件下，可食用率 $p$ 与加工时间 $t$（单位：分钟）满足函数关系 $p=at^2+bt+c$（$a$，$b$，$c$ 是常数），下图记录了三次实验的数据。根据上述函数模型和实验数据，可以得到最佳加工时间为\mbox{（　　）}
\nopagebreak[4]
\begin{center}% ============================================================
% f2014_beijing_arts_q08_popcorn.tex
% 2014 北京卷（文科）第 8 题：可食用率 p 与加工时间 t 的三次实验数据
% 来源 PDF：普通高考/2014/2014北京文.pdf
% 原页：第 1 页（A3），印面 863，第 8 题题干上方偏右
% 三个数据点：(3, 0.7)、(4, 0.8)、(5, 0.5)
% 横轴 t 每 1 分钟 1.5cm；纵轴 p 每 0.1 为 1cm。
% ============================================================
\begin{tikzpicture}[>=Stealth, line width=0.6pt, font=\small]
  % 坐标轴
  \draw[->] (0,0) -- (0,8.55) node[above] {$p$};
  \draw[->] (0,0) -- (8.8,0) node[below left=-1pt] {$t$};
  \node[below left] at (0,0) {$O$};

  % 纵轴刻度值
  \node[left] at (0,5) {$0.5$};
  \node[left] at (0,7) {$0.7$};
  \node[left] at (0,8) {$0.8$};

  % 横轴刻度值
  \node[below] at (4.5,0) {$3$};
  \node[below] at (6.0,0) {$4$};
  \node[below] at (7.5,0) {$5$};

  % 参考虚线
  \draw[dashed, line width=0.45pt] (0,7.0) -- (4.5,7.0);
  \draw[dashed, line width=0.45pt] (0,8.0) -- (6.0,8.0);
  \draw[dashed, line width=0.45pt] (0,5.0) -- (7.5,5.0);
  \draw[dashed, line width=0.45pt] (4.5,0) -- (4.5,7.0);
  \draw[dashed, line width=0.45pt] (6.0,0) -- (6.0,8.0);
  \draw[dashed, line width=0.45pt] (7.5,0) -- (7.5,5.0);

  % 三个数据点
  \fill (4.5,7.0) circle (1.5pt);
  \fill (6.0,8.0) circle (1.5pt);
  \fill (7.5,5.0) circle (1.5pt);
\end{tikzpicture}
\end{center}
\nopagebreak[4]
（A）$3.50$ 分钟；（B）$3.75$ 分钟；（C）$4.00$ 分钟；（D）$4.25$ 分钟。

\Needspace{4\baselineskip}
\textbf{解析}

\Needspace{5\baselineskip}
把三个数据点 $(3,0.7)$，$(4,0.8)$，$(5,0.5)$ 代入 $p=at^2+bt+c$：
\nopagebreak[4]
\[\begin{cases}
9a+3b+c=0.7,\\
16a+4b+c=0.8,\\
25a+5b+c=0.5 .
\end{cases}\]
\Needspace{5\baselineskip}
第二式减第一式得 $7a+b=0.1$，第三式减第二式得 $9a+b=-0.3$，两式相减得
\nopagebreak[4]
\[2a=-0.4,\quad a=-0.2,\quad b=0.1-7a=1.5,\]
\Needspace{5\baselineskip}
再由第一式 $c=0.7-9a-3b=0.7+1.8-4.5=-2$。故
\nopagebreak[4]
\[p=-0.2t^2+1.5t-2=-0.2\left(t-\frac{15}{4}\right)^2+\frac{13}{16},\]
当 $t=\dfrac{15}{4}=3.75$ 时可食用率最大，最佳加工时间为 $3.75$ 分钟，选 B。

\gkanswer{答案：B。由三点解得 $p=-0.2t^2+1.5t-2$，顶点 $t=3.75$。}

\end{answer}
\begin{answer}
\gksolhead{9}{}

\textbf{题目}：

若 $(x+\mathrm i)\mathrm i=-1+2\mathrm i$（$x\in\mathbf R$），则 $x=$\gkblank{}。

\Needspace{4\baselineskip}
\textbf{解析}

\Needspace{5\baselineskip}
展开左边，注意 $\mathrm i^2=-1$：
\nopagebreak[4]
\[(x+\mathrm i)\mathrm i=x\mathrm i+\mathrm i^2=-1+x\mathrm i .\]
与 $-1+2\mathrm i$ 比较实部、虚部，得 $x=2$。

\gkanswer{答案：$2$。$(x+\mathrm i)\mathrm i=-1+x\mathrm i=-1+2\mathrm i$，比较虚部得 $x=2$。}

\end{answer}
\begin{answer}
\gksolhead{10}{}

\textbf{题目}：

设双曲线 $C$ 的两个焦点为 $(-\sqrt2,0)$，$(\sqrt2,0)$，一个顶点是 $(1,0)$，则 $C$ 的方程为\gkblank{}。

\Needspace{4\baselineskip}
\textbf{解析}

两焦点在 $x$ 轴上且关于原点对称，故双曲线是标准形 $\dfrac{x^2}{a^2}-\dfrac{y^2}{b^2}=1$，且 $c=\sqrt2$。

\Needspace{5\baselineskip}
顶点 $(1,0)$ 在 $x$ 轴上，故 $a=1$，从而
\nopagebreak[4]
\[b^2=c^2-a^2=2-1=1 .\]
所求方程为 $x^2-y^2=1$。

\gkanswer{答案：$x^2-y^2=1$。$c=\sqrt2$，$a=1$，故 $b^2=c^2-a^2=1$。}

\end{answer}
\begin{answer}
\gksolhead{11}{}

\textbf{题目}：

某三棱锥的三视图如图所示，则该三棱锥的最长棱的棱长为\gkblank{}。
\nopagebreak[4]
\begin{center}% ============================================================
% f2014_beijing_arts_q11_threeviews.tex
% 2014 北京卷（文科）第 11 题：三棱锥的三视图
% 来源 PDF：普通高考/2014/2014北京文.pdf
% 原页：第 1 页（A3），印面 863，第 11 题题干下方
% 正（主）视图：底 2、高 2 的直角三角形，另有自顶点引向底边中点的投影线；
% 侧（左）视图：底 1、高 2 的直角三角形；俯视图：底 2、顶点在底边中点的三角形。
% ============================================================
\begin{tikzpicture}[>=Stealth, line width=0.6pt, font=\small]
  % ---------- 正（主）视图 ----------
  \draw (0,0) -- (3,0) -- (0,3) -- cycle;
  \draw (0,3) -- (1.5,0);
  % 尺寸：高 2
  \draw[<->, line width=0.45pt] (-0.55,0) -- (-0.55,3);
  \draw[line width=0.4pt] (-0.75,0) -- (-0.35,0);
  \draw[line width=0.4pt] (-0.75,3) -- (-0.35,3);
  \node[left] at (-0.60,1.5) {$2$};
  % 尺寸：底 2
  \draw[<->, line width=0.45pt] (0,-0.55) -- (3,-0.55);
  \draw[line width=0.4pt] (0,-0.75) -- (0,-0.35);
  \draw[line width=0.4pt] (3,-0.75) -- (3,-0.35);
  \node[below] at (1.5,-0.52) {$2$};
  \node at (1.5,-1.45) {正（主）视图};

  % ---------- 侧（左）视图 ----------
  \draw (6.4,0) -- (7.9,0) -- (6.4,3) -- cycle;
  \draw[<->, line width=0.45pt] (6.4,-0.55) -- (7.9,-0.55);
  \draw[line width=0.4pt] (6.4,-0.75) -- (6.4,-0.35);
  \draw[line width=0.4pt] (7.9,-0.75) -- (7.9,-0.35);
  \node[below] at (7.15,-0.52) {$1$};
  \node at (7.15,-1.45) {侧（左）视图};

  % ---------- 俯视图 ----------
  \draw (0,-3.05) -- (3,-3.05) -- (1.5,-4.55) -- cycle;
  % 尺寸：1 + 1
  \draw[<->, line width=0.45pt] (0,-2.55) -- (1.5,-2.55);
  \draw[<->, line width=0.45pt] (1.5,-2.55) -- (3,-2.55);
  \draw[line width=0.4pt] (0,-3.25) -- (0,-2.35);
  \draw[line width=0.4pt] (1.5,-3.25) -- (1.5,-2.35);
  \draw[line width=0.4pt] (3,-3.25) -- (3,-2.35);
  \node[above] at (0.75,-2.55) {$1$};
  \node[above] at (2.25,-2.55) {$1$};
  \node at (1.5,-5.05) {俯视图};
\end{tikzpicture}
\end{center}
\nopagebreak[4]
\Needspace{4\baselineskip}
\textbf{解析}

\Needspace{5\baselineskip}
由俯视图，底面三角形的一边长为 $2$，第三个顶点在该边中点的正前方；由侧视图，该顶点到底边的距离为 $1$；由正视图，底面边长 $2$、高为 $2$。取坐标
\nopagebreak[4]
\[A(0,0,0),\quad B(2,0,0),\quad C(1,1,0),\quad P(0,0,2),\]
则它在三个坐标面上的投影恰为题给三视图（正视图中的斜线是棱 $PC$ 的投影，$C$ 的投影落在 $AB$ 上）。

\Needspace{5\baselineskip}
各棱长：
\nopagebreak[4]
\[\begin{aligned}
&PA=2,\quad AB=2,\quad PB=\sqrt{2^2+2^2}=2\sqrt2,\\
&AC=BC=\sqrt2,\quad PC=\sqrt{1+1+4}=\sqrt6 .
\end{aligned}\]
最长棱为 $PB$，棱长 $2\sqrt2$。

\gkanswer{答案：$2\sqrt2$。还原得 $AB=2$、$BC$ 方向深 $1$、高 $2$，最长棱 $\sqrt{2^2+2^2}=2\sqrt2$。}

\end{answer}
\begin{answer}
\gksolhead{12}{}

\textbf{题目}：

在 $\triangle ABC$ 中，$a=1$，$b=2$，$\cos C=\dfrac14$，则 $c=$\gkblank{}；$\sin A=$\gkblank{}。

\Needspace{4\baselineskip}
\textbf{解析}

\Needspace{5\baselineskip}
由余弦定理，
\nopagebreak[4]
\[c^2=a^2+b^2-2ab\cos C=1+4-2\cdot1\cdot2\cdot\frac14=4,\]
故 $c=2$。

\Needspace{5\baselineskip}
再由余弦定理求 $\cos A$：
\nopagebreak[4]
\[\cos A=\frac{b^2+c^2-a^2}{2bc}=\frac{4+4-1}{2\cdot2\cdot2}=\frac78,\]
\Needspace{5\baselineskip}
于是
\nopagebreak[4]
\[\sin A=\sqrt{1-\cos^2A}=\sqrt{1-\frac{49}{64}}=\frac{\sqrt{15}}8 .\]
（也可由 $\sin C=\sqrt{1-\frac1{16}}=\frac{\sqrt{15}}4$ 与 $\dfrac{a}{\sin A}=\dfrac{c}{\sin C}$ 得 $\sin A=\dfrac{a\sin C}{c}=\dfrac{\sqrt{15}}8$，两法一致。）

\gkanswer{答案：$c=2$；$\sin A=\dfrac{\sqrt{15}}8$。$c^2=1+4-2\cdot1\cdot2\cdot\dfrac14=4$，$\cos A=\dfrac78$。}

\end{answer}
\begin{answer}
\gksolhead{13}{}

\textbf{题目}：

\Needspace{5\baselineskip}
若 $x$，$y$ 满足
\nopagebreak[4]
\[\begin{cases}
y\leqslant 1,\\
x-y-1\leqslant 0,\\
x+y-1\geqslant 0,
\end{cases}\]
则 $z=\sqrt3x+y$ 的最小值为\gkblank{}。

\Needspace{4\baselineskip}
\textbf{解析}

\Needspace{5\baselineskip}
约束条件即 $y\leqslant 1$，$y\geqslant x-1$，$y\geqslant 1-x$，可行域是三角形，顶点为三条边界两两的交点：
\nopagebreak[4]
\[(0,1),\qquad (2,1),\qquad (1,0).\]
\Needspace{5\baselineskip}
线性目标函数在可行域上必在顶点处取到最值，逐点计算：
\nopagebreak[4]
\[z(0,1)=1,\qquad z(2,1)=2\sqrt3+1,\qquad z(1,0)=\sqrt3 .\]
最小值为 $1$（在 $(0,1)$ 处取得）。

\gkanswer{答案：$1$。可行域顶点为 $(0,1)$，$(2,1)$，$(1,0)$，$z$ 值分别为 $1$，$2\sqrt3+1$，$\sqrt3$。}

\end{answer}
\begin{answer}
\gksolhead{14}{}

\textbf{题目}：

顾客请一位工艺师把 $A$，$B$ 两件玉石原料各制成一件工艺品，工艺师带一位徒弟完成这项任务，每件原料先由徒弟完成粗加工，再由工艺师进行精加工完成制作，两件工艺品都完成后交付顾客，两件原料每道工序所需时间（单位：工作日）如下：

\begin{center}
\begin{tabular}{lll}
\toprule
工序 & 粗加工 & 精加工 \\
\midrule
原料 $A$ & $9$ & $15$ \\
原料 $B$ & $6$ & $21$ \\
\bottomrule
\end{tabular}
\end{center}

则最短交货期为\gkblank{}个工作日。

\Needspace{4\baselineskip}
\textbf{解析}

\Needspace{5\baselineskip}
徒弟粗加工两件共需 $9+6=15$ 个工作日，工艺师精加工两件共需 $15+21=36$ 个工作日，且每件必须先粗加工后精加工，故工艺师最早只能在第 $\min\{9,6\}=6$ 个工作日开工。于是交货期
\nopagebreak[4]
\[T\geqslant 6+36=42\ \text{（\text{个工作日}）}.\]
该下界可以达到：徒弟第 $0$—$6$ 天粗加工 $B$，第 $6$—$15$ 天粗加工 $A$；工艺师第 $6$—$27$ 天精加工 $B$，第 $27$—$42$ 天精加工 $A$。两件都在第 $42$ 天完成，故最短交货期为 $42$ 个工作日。

（若先粗加工 $A$，则工艺师最早第 $9$ 天开工，交货期不早于 $45$ 天，更差。）

\gkanswer{答案：$42$。工艺师精加工共 $36$ 个工作日，最早第 $6$ 天开工，且该下界可由「先粗加工 $B$ 再粗加工 $A$」达到。}

\end{answer}
\begin{answer}
\gksolhead{15}{}

\textbf{题目}：

已知 $\{a_n\}$ 是等差数列，满足 $a_1=3$，$a_4=12$，数列 $\{b_n\}$ 满足 $b_1=4$，$b_4=20$，且 $\{b_n-a_n\}$ 是等比数列。

（1） 求数列 $\{a_n\}$ 和 $\{b_n\}$ 的通项公式；

（2） 求数列 $\{b_n\}$ 的前 $n$ 项和。

\Needspace{4\baselineskip}
\textbf{解析}

\gksubpart{（1）}

\Needspace{5\baselineskip}
设 $\{a_n\}$ 的公差为 $d$，则 $a_4=a_1+3d$，即 $12=3+3d$，$d=3$，故
\nopagebreak[4]
\[a_n=3+3(n-1)=3n .\]

\Needspace{5\baselineskip}
设 $c_n=b_n-a_n$，则 $\{c_n\}$ 是等比数列，且
\nopagebreak[4]
\[c_1=b_1-a_1=4-3=1,\qquad c_4=b_4-a_4=20-12=8 .\]
\Needspace{5\baselineskip}
设公比为 $q$，则 $q^3=\dfrac{c_4}{c_1}=8$，$q=2$，故 $c_n=2^{n-1}$，从而
\nopagebreak[4]
\[b_n=a_n+c_n=3n+2^{n-1} .\]
检验：$b_1=3+1=4$，$b_4=12+8=20$，与题设相符。

\gksubpart{（2）}

\nopagebreak[4]
\[\begin{aligned}
S_n&=\sum_{k=1}^{n}b_k=\sum_{k=1}^{n}\left(3k+2^{k-1}\right)\\
&=3\cdot\frac{n(n+1)}2+\frac{2^n-1}{2-1}\\
&=\frac{3n(n+1)}2+2^n-1 .
\end{aligned}\]
检验 $n=4$：$\dfrac{3\cdot4\cdot5}2+2^4-1=30+15=45$，而 $b_1+b_2+b_3+b_4=4+8+13+20=45$，一致。

\gkanswer{答案（1）：$a_n=3n$，$b_n=3n+2^{n-1}$。公差 $d=3$，等比数列 $\{b_n-a_n\}$ 首项 $1$、公比 $2$。}

\gkanswer{答案（2）：$S_n=\dfrac{3n(n+1)}2+2^n-1$。分组求和：等差数列部分与等比数列部分分别求和。}

\end{answer}
\begin{answer}
\gksolhead{16}{}

\textbf{题目}：

函数 $f(x)=3\sin\left(2x+\dfrac\pi6\right)$ 的部分图象如图所示。
\nopagebreak[4]
\begin{center}% ============================================================
% f2014_beijing_arts_q16_trig.tex
% 2014 北京卷（文科）第 16 题：f(x)=3sin(2x+pi/6) 的部分图象
% 来源 PDF：普通高考/2014/2014北京文.pdf
% 原页：第 1 页（A3），印面 863，右栏第 16 题题干下方
% 图中标出第二个最大值点的横坐标 x0 与纵坐标 y0。
% 横、纵同尺度：1 单位 = 1cm。
% ============================================================
\begin{tikzpicture}[>=Stealth, line width=0.8pt, font=\small]
  % 坐标轴
  \draw[->] (-0.95,0) -- (7.25,0) node[above right=-2pt] {$x$};
  \draw[->] (0,-3.75) -- (0,4.05) node[left] {$y$};
  \node[below right] at (0.02,-0.05) {$O$};

  % 图象
  \draw[domain=-0.5:6.6, samples=240, smooth]
    plot (\x,{3*sin(deg(2*\x+pi/6))});

  % x0、y0 的参考虚线
  \draw[dashed, line width=0.45pt] (0,3) -- (3.665,3);
  \draw[dashed, line width=0.45pt] (3.665,0) -- (3.665,3);
  \node[left] at (-0.05,3) {$y_0$};
  \node[below] at (3.665,0) {$x_0$};
\end{tikzpicture}
\end{center}
\nopagebreak[4]
（1） 写出 $f(x)$ 的最小正周期及图中 $x_0$，$y_0$ 的值；

（2） 求 $f(x)$ 在区间 $\left[-\dfrac\pi2,-\dfrac\pi{12}\right]$ 上的最大值和最小值。

\Needspace{4\baselineskip}
\textbf{解析}

\gksubpart{（1）}

\Needspace{5\baselineskip}
最小正周期
\nopagebreak[4]
\[T=\frac{2\pi}{2}=\pi .\]
\Needspace{5\baselineskip}
$f(x)$ 取最大值 $3$ 当且仅当
\nopagebreak[4]
\[2x+\frac\pi6=\frac\pi2+2k\pi,\quad\text{\text{即}}\quad x=\frac\pi6+k\pi\ (k\in\mathbf Z).\]
\Needspace{5\baselineskip}
图象中 $x_0$ 是 $y$ 轴右侧第二个最大值点的横坐标，取 $k=1$ 得
\nopagebreak[4]
\[x_0=\frac\pi6+\pi=\frac{7\pi}6,\qquad y_0=3 .\]

\gksubpart{（2）}

\Needspace{5\baselineskip}
令 $u=2x+\dfrac\pi6$。当 $x\in\left[-\dfrac\pi2,-\dfrac\pi{12}\right]$ 时，
\nopagebreak[4]
\[u\in\left[-\frac{5\pi}6,\,0\right].\]
\Needspace{5\baselineskip}
$\sin u$ 在 $\left[-\dfrac{5\pi}6,-\dfrac\pi2\right]$ 上递减、在 $\left[-\dfrac\pi2,0\right]$ 上递增，故最小值在 $u=-\dfrac\pi2$ 处：
\nopagebreak[4]
\[\sin\left(-\frac\pi2\right)=-1,\qquad f_{\min}=-3,\]
此时 $x=\dfrac12\left(-\dfrac\pi2-\dfrac\pi6\right)=-\dfrac\pi3\in\left[-\dfrac\pi2,-\dfrac\pi{12}\right]$。

\Needspace{5\baselineskip}
最大值在端点处取得，比较
\nopagebreak[4]
\[\sin\left(-\frac{5\pi}6\right)=-\frac12,\qquad \sin 0=0,\]
故最大值为 $3\sin0=0$，在 $x=-\dfrac\pi{12}$ 处取得。

所以最大值为 $0$，最小值为 $-3$。

\gkanswer{答案（1）：$T=\pi$，$x_0=\dfrac{7\pi}6$，$y_0=3$。最大值点 $x=\dfrac\pi6+k\pi$，图中 $x_0$ 为第二个最大值点。}

\gkanswer{答案（2）：最大值为 $0$，最小值为 $-3$。$2x+\dfrac\pi6\in\left[-\dfrac{5\pi}6,0\right]$，$\sin$ 在 $u=-\dfrac\pi2$ 处取 $-1$，端点处最大为 $\sin0=0$。}

\end{answer}
\begin{answer}
\gksolhead{17}{}

\textbf{题目}：

如图，在三棱柱 $ABC\text{-}A_1B_1C_1$ 中，侧棱垂直于底面，$AB\perp BC$，$AA_1=AC=2$，$BC=1$，$E$，$F$ 分别为 $A_1C_1$，$BC$ 的中点。
\nopagebreak[4]
\begin{center}% ============================================================
% f2014_beijing_arts_q17_prism.tex
% 2014 北京卷（文科）第 17 题：直三棱柱 ABC-A1B1C1
% 来源 PDF：普通高考/2014/2014北京文.pdf
% 原页：第 1 页（A3），印面 863，右栏第 17 题题干下方
% AB 垂直 BC，AA1=AC=2，BC=1；E 为 A1C1 中点，F 为 BC 中点。
% 图中另画出 A1B、EB、B1C（虚线）与 C1F（实线）。
% ============================================================
\begin{tikzpicture}[>=Stealth, line width=0.65pt, font=\small,
  line cap=round, line join=round]
  % 顶点
  \coordinate (A1) at (0,4.50);
  \coordinate (B1) at (1.75,3.65);
  \coordinate (C1) at (3.40,4.15);
  \coordinate (A)  at (0,1.00);
  \coordinate (B)  at (1.75,0.15);
  \coordinate (C)  at (3.40,0.65);
  \coordinate (E)  at (1.70,4.325); % A1C1 中点
  \coordinate (F)  at (2.575,0.40); % BC 中点

  % 被遮挡的棱与辅助线
  \draw[dashed] (A) -- (C);
  \draw[dashed] (A1) -- (B);
  \draw[dashed] (B1) -- (C);
  \draw[dashed] (E) -- (B);

  % 可见的棱
  \draw (A1) -- (B1) -- (C1);
  \draw (A1) -- (C1);
  \draw (A1) -- (A);
  \draw (B1) -- (B);
  \draw (C1) -- (C);
  \draw (A) -- (B) -- (C);
  \draw (C1) -- (F);

  % 顶点标注
  \node[above left=-1pt] at (A1) {$A_1$};
  \node[above] at (E) {$E$};
  \node[above right=-1pt] at (C1) {$C_1$};
  \node[right] at (B1) {$B_1$};
  \node[below left=-1pt] at (A) {$A$};
  \node[below] at (B) {$B$};
  \node[right] at (C) {$C$};
  \node[below right=-1pt] at (F) {$F$};
\end{tikzpicture}
\end{center}
\nopagebreak[4]
（1） 求证：平面 $ABE\perp$ 平面 $B_1BCC_1$；

（2） 求证：$C_1F\parallel$ 平面 $ABE$；

（3） 求三棱锥 $E\text{-}ABC$ 的体积。

\Needspace{4\baselineskip}
\textbf{解析}

\gksubpart{（1）}

\Needspace{5\baselineskip}
侧棱垂直于底面，故 $BB_1\perp$ 平面 $ABC$，而 $AB\subset$ 平面 $ABC$，所以
\nopagebreak[4]
\[BB_1\perp AB .\]
\Needspace{5\baselineskip}
又已知 $AB\perp BC$，且 $BB_1\cap BC=B$，$BB_1$，$BC\subset$ 平面 $B_1BCC_1$，由线面垂直判定定理得
\nopagebreak[4]
\[AB\perp \text{\text{平面} }B_1BCC_1 .\]
而 $AB\subset$ 平面 $ABE$，由面面垂直判定定理，平面 $ABE\perp$ 平面 $B_1BCC_1$。

\gksubpart{（2）}

取 $AB$ 的中点 $G$，连接 $EG$，$FG$。

\Needspace{5\baselineskip}
在 $\triangle ABC$ 中，$F$，$G$ 分别为 $BC$，$AB$ 的中点，故
\nopagebreak[4]
\[FG\parallel AC,\qquad FG=\frac12AC=1 .\]
\Needspace{5\baselineskip}
在三棱柱中 $A_1C_1\parallel AC$ 且 $A_1C_1=AC=2$；$E$ 为 $A_1C_1$ 的中点，故
\nopagebreak[4]
\[EC_1\parallel AC,\qquad EC_1=\frac12A_1C_1=1 .\]
\Needspace{5\baselineskip}
于是 $FG\parallel EC_1$ 且 $FG=EC_1$，四边形 $FGEC_1$ 是平行四边形，从而
\nopagebreak[4]
\[C_1F\parallel EG .\]
又 $EG\subset$ 平面 $ABE$，$C_1F\not\subset$ 平面 $ABE$，由线面平行判定定理得 $C_1F\parallel$ 平面 $ABE$。

\gksubpart{（3）}

\Needspace{5\baselineskip}
底面 $\triangle ABC$ 中 $AB\perp BC$，$AC=2$，$BC=1$，故
\nopagebreak[4]
\[AB=\sqrt{AC^2-BC^2}=\sqrt3,\qquad
S_{\triangle ABC}=\frac12\cdot AB\cdot BC=\frac{\sqrt3}2 .\]
\Needspace{5\baselineskip}
$E$ 在上底面 $A_1B_1C_1$ 上，而上下底面平行、侧棱垂直于底面，故 $E$ 到平面 $ABC$ 的距离等于侧棱长
\nopagebreak[4]
\[h=AA_1=2 .\]
\Needspace{5\baselineskip}
所以
\nopagebreak[4]
\[V_{E\text{-}ABC}=\frac13S_{\triangle ABC}\cdot h=\frac13\cdot\frac{\sqrt3}2\cdot2=\frac{\sqrt3}3 .\]

\gkanswer{答案（1）：见解析版}

\gkanswer{答案（2）：见解析版}

\gkanswer{答案（3）：$\dfrac{\sqrt3}3$。$AB=\sqrt3$，$S_{\triangle ABC}=\dfrac{\sqrt3}2$，高为 $AA_1=2$，$V=\dfrac13\cdot\dfrac{\sqrt3}2\cdot2$。}

\end{answer}
\begin{answer}
\gksolhead{18}{}

\textbf{题目}：

从某校随机抽取 $100$ 名学生，获得了他们一周课外阅读时间（单位：小时）的数据，整理得到数据分组及频数分布表和频率分布直方图：

\begin{center}
\begin{tabular}{lll}
\toprule
组号 & 分组 & 频数 \\
\midrule
$1$ & $[0,2)$ & $6$ \\
$2$ & $[2,4)$ & $8$ \\
$3$ & $[4,6)$ & $17$ \\
$4$ & $[6,8)$ & $22$ \\
$5$ & $[8,10)$ & $25$ \\
$6$ & $[10,12)$ & $12$ \\
$7$ & $[12,14)$ & $6$ \\
$8$ & $[14,16)$ & $2$ \\
$9$ & $[16,18)$ & $2$ \\
合计 &  & $100$ \\
\bottomrule
\end{tabular}
\end{center}
\nopagebreak[4]
\begin{center}% ============================================================
% f2014_beijing_arts_q18_histogram.tex
% 2014 北京卷（文科）第 18 题：100 名学生一周课外阅读时间的频率分布直方图
% 来源 PDF：普通高考/2014/2014北京文.pdf
% 原页：第 2 页（A3），印面 864，第 18 题频数分布表右侧
% 组距 2，共 9 组；柱高 = 频数/200，其中第 3 组顶端标 a、第 5 组顶端标 b。
% 横轴每 2 小时 1.5cm，纵轴每 0.025 为 0.75cm（b=0.125 对应 3.75cm）。
% ============================================================
\begin{tikzpicture}[>=Stealth, line width=0.6pt, font=\small,
  x=0.75cm, y=0.75cm]
  % 坐标轴
  \draw[->] (0,0) -- (0,5.75);
  \draw[->] (0,0) -- (20.6,0);

  % 纵轴标题（频率/组距）
  \draw (0.05,6.40) -- (1.55,6.40);
  \node[above] at (0.80,6.40) {频率};
  \node[below] at (0.80,6.40) {组距};

  % 横轴刻度
  \node[below left] at (0,0) {$0$};
  \foreach \x in {2,4,6,8,10,12,14,16,18} {
    \draw (\x,0) -- (\x,-0.14);
    \node[below] at (\x,-0.02) {$\x$};
  }

  % 九个矩形：频数 6,8,17,22,25,12,6,2,2
  \foreach \i/\n in {0/6,1/8,2/17,3/22,4/25,5/12,6/6,7/2,8/2} {
    \draw (\i*2,0) rectangle (\i*2+2,\n*0.16);
  }

  % a、b 的参考虚线与标注
  \draw[dashed, line width=0.45pt] (0,2.72) -- (6,2.72);
  \draw[dashed, line width=0.45pt] (0,4.00) -- (10,4.00);
  \node[left] at (0,2.72) {$a$};
  \node[left] at (0,4.00) {$b$};

  % 横轴名称
  \node[above] at (19.1,0.12) {阅读时间};
\end{tikzpicture}
\end{center}
\nopagebreak[4]
（1） 从该校随机选取一名学生，试估计这名学生该周课外阅读时间少于 $12$ 小时的概率；

（2） 求频率分布直方图中的 $a$，$b$ 的值；

（3） 假设同一组中的每个数据可用该组区间的中点值代替，试估计样本中的 $100$ 名学生该周课外阅读时间的平均数在第几组。（只需写出结论）

\Needspace{4\baselineskip}
\textbf{解析}

\gksubpart{（1）}

\Needspace{5\baselineskip}
阅读时间少于 $12$ 小时即落在前 $6$ 组，频数为
\nopagebreak[4]
\[6+8+17+22+25+12=90,\]
\Needspace{5\baselineskip}
故概率的估计值为
\nopagebreak[4]
\[\frac{90}{100}=0.9 .\]

\gksubpart{（2）}

\Needspace{5\baselineskip}
组距为 $2$，直方图中矩形的高为 $\dfrac{\text{\text{频数}}}{100\times2}=\dfrac{\text{\text{频数}}}{200}$。第 $3$ 组频数 $17$，第 $5$ 组频数 $25$，故
\nopagebreak[4]
\[a=\frac{17}{200}=0.085,\qquad b=\frac{25}{200}=0.125 .\]

\gksubpart{（3）}

\Needspace{5\baselineskip}
用各组中点值代表该组数据，中点依次为 $1,3,5,7,9,11,13,15,17$，平均数
\nopagebreak[4]
\[\begin{aligned}
\bar x&=\frac{1}{100}\bigl(1\times6+3\times8+5\times17+7\times22+9\times25\\
&\qquad\qquad +11\times12+13\times6+15\times2+17\times2\bigr)\\
&=\frac{6+24+85+154+225+132+78+30+34}{100}=\frac{768}{100}=7.68 .
\end{aligned}\]
$7.68\in[6,8)$，故平均数在第 $4$ 组。

\gkanswer{答案（1）：$0.9$。少于 $12$ 小时的频数为 $90$，$\dfrac{90}{100}=0.9$。}

\gkanswer{答案（2）：$a=0.085$，$b=0.125$。矩形高为 $\dfrac{\text{\text{频数}}}{200}$，第 $3$ 组 $17$、第 $5$ 组 $25$。}

\gkanswer{答案（3）：第 $4$ 组。中点值加权平均得 $\bar x=7.68\in[6,8)$。}

\end{answer}
\begin{answer}
\gksolhead{19}{}

\textbf{题目}：

已知椭圆 $C:x^2+2y^2=4$。

（1） 求椭圆 $C$ 的离心率；

（2） 设 $O$ 为原点，若点 $A$ 在直线 $y=2$ 上，点 $B$ 在椭圆 $C$ 上，且 $OA\perp OB$，求线段 $AB$ 长度的最小值。

\Needspace{4\baselineskip}
\textbf{解析}

\gksubpart{（1）}

\Needspace{5\baselineskip}
把方程写成标准形
\nopagebreak[4]
\[\frac{x^2}4+\frac{y^2}2=1,\]
\Needspace{5\baselineskip}
故 $a^2=4$，$b^2=2$，$c^2=a^2-b^2=2$，于是
\nopagebreak[4]
\[e=\frac ca=\frac{\sqrt2}{2} .\]

\gksubpart{（2）}

\Needspace{5\baselineskip}
设 $A(s,2)$，$B(x,y)$，其中 $x^2+2y^2=4$。由 $OA\perp OB$ 得
\nopagebreak[4]
\[\overrightarrow{OA}\cdot\overrightarrow{OB}=sx+2y=0 .\]
因为 $B$ 在椭圆上，$|x|\leqslant2$，而 $A$ 在 $y=2$ 上，故 $y\neq2$，$B\neq A$；当 $x\neq0$ 时 $s=-\dfrac{2y}{x}$，当 $x=0$ 时 $y=\pm\sqrt2$，由 $sx+2y=0$ 得 $y=0$，矛盾，故只需考虑 $x\neq0$。

\Needspace{5\baselineskip}
由于 $\angle AOB=90^\circ$，
\nopagebreak[4]
\[|AB|^2=|OA|^2+|OB|^2 .\]
\Needspace{5\baselineskip}
由 $x^2=4-2y^2$ 得
\nopagebreak[4]
\[|OA|^2=s^2+4=\frac{4y^2}{x^2}+4=\frac{4y^2}{4-2y^2}+4
=\frac{2y^2}{2-y^2}+4=\frac{2(4-y^2)}{2-y^2},\]
\Needspace{5\baselineskip}
又
\nopagebreak[4]
\[|OB|^2=x^2+y^2=4-y^2 .\]
\Needspace{5\baselineskip}
令 $u=2-y^2$，由 $y^2\in[0,2)$ 知 $u\in(0,2]$，且 $4-y^2=u+2$，于是
\nopagebreak[4]
\[|AB|^2=\frac{2(u+2)}{u}+(u+2)=\frac{(u+2)^2}{u}=u+4+\frac4u\geqslant 4+2\sqrt{u\cdot\frac4u}=8,\]
等号当且仅当 $u=2$，即 $y=0$，$x=\pm2$，此时 $s=0$，$A(0,2)$，$B(\pm2,0)$，确实满足 $OA\perp OB$。

所以 $|AB|$ 的最小值为 $\sqrt8=2\sqrt2$。

\gkanswer{答案（1）：$\dfrac{\sqrt2}2$。$a=2$，$b=\sqrt2$，$c=\sqrt2$。}

\gkanswer{答案（2）：$2\sqrt2$。$|AB|^2=|OA|^2+|OB|^2=\dfrac{(u+2)^2}{u}\geqslant 8$，其中 $u=2-y^2\in(0,2]$，$u=2$ 时取等号。}

\end{answer}
\begin{answer}
\gksolhead{20}{}

\textbf{题目}：

已知函数 $f(x)=2x^3-3x$。

（1） 求 $f(x)$ 在区间 $[-2,1]$ 上的最大值；

（2） 若过点 $P(1,t)$ 存在 $3$ 条直线与曲线 $y=f(x)$ 相切，求 $t$ 的取值范围；

（3） 问过点 $A(-1,2)$，$B(2,10)$，$C(0,2)$ 分别存在几条直线与曲线 $y=f(x)$ 相切？（只需写出结论）

\Needspace{4\baselineskip}
\textbf{解析}

\gksubpart{（1）}

\nopagebreak[4]
\[f'(x)=6x^2-3=3\left(2x^2-1\right),\]
零点为 $x=\pm\dfrac{\sqrt2}2$。在 $[-2,1]$ 上列表：

\begin{center}
\begingroup\setlength{\tabcolsep}{4pt}
\begin{tabular}{llllllll}
\toprule
$x$ & $-2$ & $\left(-2,-\frac{\sqrt2}2\right)$ & $-\frac{\sqrt2}2$ & $\left(-\frac{\sqrt2}2,\frac{\sqrt2}2\right)$ & $\frac{\sqrt2}2$ & $\left(\frac{\sqrt2}2,1\right)$ & $1$ \\
\midrule
$f'(x)$ &  & $+$ & $0$ & $-$ & $0$ & $+$ &  \\
$f(x)$ & $-10$ & 增 & $\sqrt2$ & 减 & $-\sqrt2$ & 增 & $-1$ \\
\bottomrule
\end{tabular}
\endgroup
\end{center}

\Needspace{5\baselineskip}
其中
\nopagebreak[4]
\[f(-2)=-16+6=-10,\quad f\left(-\frac{\sqrt2}2\right)=\frac{\sqrt2}2+\frac{3\sqrt2}2=\sqrt2,\quad f\left(\frac{\sqrt2}2\right)=-\sqrt2,\quad f(1)=-1 .\]
比较得 $f(x)$ 在 $[-2,1]$ 上的最大值为 $\sqrt2$。

\gksubpart{（2）}

\Needspace{5\baselineskip}
设切点为 $\left(x_0,\,2x_0^3-3x_0\right)$，则切线斜率为 $6x_0^2-3$，切线方程为
\nopagebreak[4]
\[y-\left(2x_0^3-3x_0\right)=\left(6x_0^2-3\right)(x-x_0).\]
\Needspace{5\baselineskip}
代入 $P(1,t)$：
\nopagebreak[4]
\[\begin{aligned}
t&=\left(6x_0^2-3\right)(1-x_0)+2x_0^3-3x_0\\
&=6x_0^2-3-6x_0^3+3x_0+2x_0^3-3x_0\\
&=-4x_0^3+6x_0^2-3 .
\end{aligned}\]
\Needspace{5\baselineskip}
故过 $P$ 的切线条数等于方程 $g(x)=-4x^3+6x^2-3=t$ 的不同实根个数，其中
\nopagebreak[4]
\[g'(x)=-12x^2+12x=12x(1-x).\]
\Needspace{5\baselineskip}
$g$ 在 $(-\infty,0)$ 上递减，在 $(0,1)$ 上递增，在 $(1,+\infty)$ 上递减，故
\nopagebreak[4]
\[g(0)=-3\ \text{\text{为极小值}},\qquad g(1)=-4+6-3=-1\ \text{\text{为极大值}},\]
\Needspace{5\baselineskip}
且 $x\to-\infty$ 时 $g\to+\infty$，$x\to+\infty$ 时 $g\to-\infty$。于是水平线 $y=t$ 与曲线 $y=g(x)$ 有三个不同交点当且仅当
\nopagebreak[4]
\[-3<t<-1 .\]
即 $t$ 的取值范围为 $(-3,-1)$。

\gksubpart{（3）}

\Needspace{5\baselineskip}
设过点 $(a,b)$ 的切线的切点横坐标为 $x_0$，同（2）的推导得
\nopagebreak[4]
\[b=-4x_0^3+6ax_0^2-3a,\quad\text{\text{即}}\quad 4x_0^3-6ax_0^2+3a+b=0 .\]
切线条数等于该三次方程不同实根的个数。

\begin{itemize}
\item $A(-1,2)$：$a=-1$，$b=2$，方程化为
\end{itemize}
\nopagebreak[4]
\[4x^3+6x^2-1=0 .\]
记 $h(x)=4x^3+6x^2-1$，$h'(x)=12x(x+1)$，$h$ 在 $x=-1$ 取极大值 $h(-1)=1>0$，在 $x=0$ 取极小值 $h(0)=-1<0$，且 $h(-\infty)=-\infty$，$h(+\infty)=+\infty$，故有三个不同实根（分别在 $(-\infty,-1)$，$(-1,0)$，$(0,+\infty)$ 内），过 $A$ 有 $3$ 条切线。

\begin{itemize}
\item $B(2,10)$：$a=2$，$b=10$，方程化为
\end{itemize}
\nopagebreak[4]
\[4x^3-12x^2+16=0\iff x^3-3x^2+4=0\iff(x+1)(x-2)^2=0,\]
有两个不同实根 $x=-1$（单根）与 $x=2$（二重根），故过 $B$ 有 $2$ 条切线（其中一条的切点就是 $B$ 本身）。

\begin{itemize}
\item $C(0,2)$：$a=0$，$b=2$，方程化为
\end{itemize}
\nopagebreak[4]
\[4x^3+2=0\iff x^3=-\frac12,\]
在实数范围内只有一个根，故过 $C$ 有 $1$ 条切线。

结论：过 $A$，$B$，$C$ 分别有 $3$ 条、$2$ 条、$1$ 条直线与曲线 $y=f(x)$ 相切。

\gkanswer{答案（1）：$\sqrt2$。$f'(x)=3(2x^2-1)$，极大值 $f\left(-\dfrac{\sqrt2}2\right)=\sqrt2$ 大于端点值 $f(-2)=-10$、$f(1)=-1$。}

\gkanswer{答案（2）：$(-3,-1)$。切点横坐标满足 $-4x^3+6x^2-3=t$，该函数极小值 $-3$、极大值 $-1$，三根当且仅当 $-3<t<-1$。}

\gkanswer{答案（3）：过 $A$ 有 $3$ 条，过 $B$ 有 $2$ 条，过 $C$ 有 $1$ 条。对应方程 $4x^3+6x^2-1=0$（三个实根）、$(x+1)(x-2)^2=0$、$4x^3+2=0$（一个实根）。}

\end{answer}